\chapter{Fibrodysplasia}
\index{Fibrodysplasia}

\section*{Definition and Overview}
Fibrodysplasia ossificans progressiva (FOP) is an extremely rare genetic disorder in which bone forms inside muscles, tendons, ligaments, and other connective tissues, progressively immobilising the body. The name means "progressive formation of bone in fibrous tissue". A person with FOP is born with malformed great toes, and in childhood begins to develop swellings that turn into bone, often triggered by injury, injections, or viral illness, but sometimes appearing without any trigger. Over time, the extra bone fuses joints, restricts movement, and eventually confines the person to immobility. There is no effective treatment, and surgery to remove the extra bone makes it worse, so management focuses on prevention, symptom control, and quality of life.

\section*{Epidemiology}
FOP is among the rarest diseases known, affecting approximately 1 in 2 million people worldwide, which translates to around one person per several hundred million. Fewer than 1,000 cases have been described globally. It affects all ethnic groups and both sexes equally, and most cases arise from a new spontaneous mutation, so there is usually no family history. Because of its rarity, the diagnosis is frequently delayed or mistaken for other conditions such as aggressive juvenile fibromatosis, and many people are harmed by biopsies or surgery before the true diagnosis is recognised. In many countries, only a handful of people live with the condition, supported by specialist teams and patient organisations.

\section*{Aetiology and Pathophysiology}
\begin{itemize}
    \item \textbf{Genetic cause:} a mutation in the ACVR1 gene, found in essentially all people with the classic form of FOP
    \item \textbf{Mechanism:} the mutation makes a receptor for bone morphogenetic protein (BMP) overactive, so cells that should repair soft tissue instead produce bone
    \item \textbf{Heterotopic ossification:} bone forms in abnormal places, including muscles, tendons, and fascia, through a process resembling normal bone formation
    \item \textbf{Triggers:} flare-ups are often provoked by trauma, intramuscular injections, immunisations, surgery, or viral illness, but can occur spontaneously
    \item \textbf{Malformed great toes:} a congenital sign present at birth, with short, bent, or single-jointed big toes
    \item \textbf{Progression:} ossification typically begins in childhood in the neck, shoulders, and spine, and spreads over decades, eventually causing near-total immobility
    \item \textbf{The difficulty of treatment:} attempts to remove the extra bone stimulate more bone formation, so surgery is contraindicated
\end{itemize}

\section*{Clinical Features}
\begin{itemize}
    \item Malformed great toes at birth, the characteristic early sign
    \item Swellings on the neck, back, or shoulders in childhood that may be warm and painful, then harden into bone over weeks to months
    \item Progressive stiffening and loss of movement of the neck, jaw, shoulders, spine, and limbs
    \item Difficulty with chewing, speaking, and breathing as the jaw and chest become involved
    \item Scoliosis and fusion of the spine
    \item Flare-ups that worsen after injury, injections, surgery, or illness
    \item Preservation of intelligence: cognitive development is normal
    \item Immobility, which is profound in later decades, with many people becoming confined to a wheelchair or bed
\end{itemize}

\section*{Differential Diagnosis}
Because of its rarity, FOP is often initially misdiagnosed.
\begin{itemize}
    \item Aggressive juvenile fibromatosis, a tumour-like overgrowth of connective tissue
    \item Tumoural calcinosis and other causes of soft-tissue calcification
    \item Progressive osseous heteroplasia, a related but distinct bone-forming disorder
    \item Fibrous dysplasia, which affects bone itself rather than soft tissue
    \item Juvenile arthritis or soft-tissue tumours when swellings first appear
    \item The malformed toes help distinguish FOP from these conditions at birth
\end{itemize}

\section*{When to Seek Medical Care}
Because surgical and procedural interventions trigger new bone formation, it is essential that every healthcare provider caring for a person with FOP is aware of the diagnosis before any procedure. Families should seek care from a specialist FOP team, and should carry information about the condition to all appointments.

\textbf{Contact your local emergency services for:}
\begin{itemize}
    \item Difficulty breathing, which can develop as chest wall movement is restricted
    \item Falls with head injury, since even minor trauma can trigger flare-ups and bleeding
    \item Severe pain with a flare-up that is not controlled
    \item Any emergency, in which it is vital to tell paramedics about the diagnosis and the need to avoid intramuscular injections
\end{itemize}

\section*{Diagnosis and Investigations}
\begin{itemize}
    \item \textbf{Clinical recognition:} the combination of malformed great toes and progressive soft-tissue ossification is almost diagnostic
    \item \textbf{Genetic testing:} identifies the ACVR1 mutation and confirms the diagnosis
    \item \textbf{Imaging:} X-rays, CT, and bone scans show the extra bone, but imaging should be used judiciously because even minor trauma can trigger flare-ups
    \item \textbf{Avoidance of biopsy:} biopsy of a swelling can itself trigger bone formation and is avoided once FOP is suspected
    \item \textbf{Audiometry and dental review:} to monitor the hearing loss and jaw problems that can develop
    \item \textbf{Respiratory assessment:} lung function is monitored as chest mobility decreases
\end{itemize}

\section*{Management}
\begin{itemize}
    \item \textbf{Prevention of triggers:} the cornerstone of care; avoid intramuscular injections, biopsies, surgery on the extra bone, and injury where possible; dental work and immunisations are planned with the specialist team
    \item \textbf{Medicines for flare-ups:} corticosteroids and other anti-inflammatory agents may reduce the swelling of acute flare-ups, though no medicine reliably stops bone formation
    \item \textbf{Symptom control:} analgesia for pain, physiotherapy within pain-free ranges to maintain movement, and splinting for comfort
    \item \textbf{Assistive technology:} equipment for mobility, positioning, and daily care as function declines
    \item \textbf{Respiratory care:} monitoring, chest physiotherapy, and support for breathing difficulties
    \item \textbf{Clinical trials:} investigational drugs targeting the BMP pathway are in development, and specialist centres can advise on eligibility
    \item \textbf{Psychosocial support:} counselling, patient organisations, and support networks for people and families
\end{itemize}

\section*{Complications}
\begin{itemize}
    \item Progressive fusion of all joints, leading to complete immobility
    \item Restriction of the chest, causing respiratory failure, the most common cause of death
    \item Difficulty eating and maintaining nutrition as the jaw fuses
    \item Hearing loss and dental problems
    \item Falls and fractures, which worsen ossification
    \item Severe pain and reduced quality of life without good supportive care
\end{itemize}

\section*{Prognosis and Outcomes}
FOP is a severe, progressive condition with no cure, and most affected people become profoundly immobile over the second to fourth decades of life. The most common cause of death is respiratory failure from chest wall restriction, with median survival into middle age, though outcomes vary. Research has transformed understanding of the condition, and clinical trials of targeted therapies are underway. For people living with FOP, expert multidisciplinary care, meticulous avoidance of triggers, and strong psychosocial support markedly improve quality of life even though they cannot halt the disease.

\section*{Prevention and Self-Care}
\begin{itemize}
    \item Educate all healthcare providers about the diagnosis and the rule against intramuscular injections, biopsies, and surgery on ossified tissue
    \item Carry a medical alert card or bracelet and information about the condition
    \item Plan dental work, immunisations, and any procedure with the specialist FOP team
    \item Avoid activities with a high risk of injury and falls
    \item Maintain gentle, pain-free movement where possible
    \item Access disability supports and assistive equipment early
    \item Connect with patient organisations for information, trials, and community support
\end{itemize}

\section*{Sources and Further Reading}
The following reputable sources provide further information for healthcare students and patients seeking reliable, current advice about fibrodysplasia ossificans progressiva.
\begin{itemize}
    \item International Fibrodysplasia Ossificans Progressiva Association. \textit{FOP resources and research.} https://www.ifopa.org/
    \item Genetic and Rare Diseases Information Center. \textit{FOP.} https://rarediseases.info.nih.gov/
    \item The Royal Children's Hospital Melbourne. \textit{Genetic services.} https://www.rch.org.au/
    \item Orphanet. \textit{Fibrodysplasia ossificans progressiva.} https://www.orpha.net/
    \item Mayo Clinic. \textit{FOP.} https://www.mayoclinic.org/
    \item MedlinePlus. \textit{Fibrodysplasia ossificans progressiva.} https://medlineplus.gov/
\end{itemize}
